\section{Conclusion}

This research successfully designed, implemented, and evaluated a high-performance web virtualization framework for Augmented Reality, directly addressing the general objective of minimizing computational overhead in browser-based AR. The TapTapp AR V11 architecture, characterized by its Nanite-style vision codec and logic-only JavaScript pipeline, reduced compilation times by over 90\% and payload sizes by 86\% compared to the industry standard.

The scientific contribution of this work lies in the validation of Fourier positional encoding and dynamic scale filtering as viable alternatives to deep convolutional networks for planar image tracking. We have demonstrated that for specific computer vision tasks, optimized algebraic algorithms can outperform generalized neural approaches by orders of magnitude on resource-constrained devices.

These optimizations are scalable and immediately applicable to the development of inclusive web technologies, ensuring that the immersive web is not restricted to users with high-end hardware. The code and protocols presented here open new avenues for "Main Thread AI," proving that JavaScript is capable of real-time, high-frequency signal processing when freed from the overhead of generic tensor engines.
